%======================================================================
\section{Modo rápido con memoria adaptable (v8)}\label{sec:rapido}

Las secciones anteriores tratan \SotSort con memoria constante, que en la
implementación de v8 sigue siendo \texttt{stable\_sort}. La versión 8 agrega dos
entradas que usan memoria adicional si la hay: \Rapido (\texttt{stable\_sort\_fast}) y
\PorClave (\texttt{stable\_sort\_by\_key}). Esta sección da su pseudocódigo y las
demostraciones, al mismo nivel de detalle que las anteriores. $\Bsz_0$ denota el buffer
propio de \SotSort (el $\Bsz$ de las secciones anteriores) y $\beta_m$ el presupuesto
de memoria en bytes.

\subsection{Supuestos adicionales y tamaño del buffer}

\begin{supuesto}[Hilos]\label{sup:hilos}
Si $P>1$, $\lt$ puede evaluarse desde varios hilos a la vez, y cada evaluación sigue
cumpliendo el Supuesto~\ref{sup:orden}. Lo que un hilo escribe antes de una barrera (la
espera \texttt{join} de C++) es visible para todos después de ella.
\end{supuesto}
\begin{supuesto}[Enteros]\label{sup:enteros}
Si $T$ es un tipo entero y $\lt$ es $<$, el estado observable de un elemento es su
valor: dos elementos con el mismo valor son indistinguibles. Para \Conteo, ``$A=\Omega$''
significa que $A$ y $\Omega$ tienen la misma secuencia de valores.
\end{supuesto}
\begin{supuesto}[Reservas]\label{sup:reserva}
\Reservar{$x, E$} o falla sin efectos, o entrega $E$ celdas disjuntas de $A$ y de toda
otra memoria en uso. En Linux con \emph{overcommit}, una reserva puede salir bien y
fallar recién al usarse; ese caso queda fuera del modelo.
\end{supuesto}

\begin{observacion}[Buffer más grande que $\Bsz_0$]\label{obs:bsz}
Revisé cada aparición de $\Bsz$ en las demostraciones anteriores. Solo se usan cotas
inferiores ($\Bsz\ge 16$ en la muestra del pivote y en la pila de \Part) y
precondiciones de la forma ``el segmento entra en el buffer'', que con un $\Bsz$ mayor se
cumplen igual o con más razón. La cota superior $\Bsz\le4096$ del
Supuesto~\ref{sup:n} describía la implementación de v7 y no interviene en ninguna
prueba. Por lo tanto los lemas y teoremas anteriores valen para todo $\Bsz\ge 16$, en
particular para $\Bsz=E$. La hoja sigue siendo $\leaf=\min(128,\Bsz_0)\le\Bsz$.
\end{observacion}

% Algoritmos del modo rápido (v8). Se incluye desde sotsort.tex y sotsort_demo.tex.
% Requiere: \lt, \nlt, \Mrg, \buf, \Bsz y las funciones \SotSort, \Ordenado,
% \PocosGiros, \Aleatorio, \MIzq, \MDer, \PlanF, \Rapido, \Galope, \Conteo,
% \CoRango, \Paralelo, \PorClave, \Reservar, \Disp.

\begin{algorithm}[H]
\caption{Modo rápido (v8): plan de memoria y punto de entrada}\label{alg:plan}
\KwIn{$A[0..n)$, comparador $\lt$, opciones $O$: tope $O.\mathit{max}$ en bytes, fracción $O.f$ ($0{,}5$), $O.\mathit{sondear}$, $O.\mathit{hilos}$ ($0$ = automático), $O.\mathit{atajos}$}
\Fn(\tcp*[h]{$s=\mathit{sizeof}(T)$; $\Bsz_0$: buffer propio, unos 16\,KB}){\PlanF{$n, s, O$}}{
  $\beta_m\gets O.\mathit{max}$\;
  \If(\tcp*[f]{solo si el buffer pasaría de 1\,MB}){$ns>2^{20}$ \KwAnd $O.\mathit{sondear}$}{
    $D\gets$ \Disp{}\tcp*{$\min(\text{MemAvailable},\ \text{margen del cgroup})$; $0$ si no se sabe}
    $\beta_m\gets\min\big(\beta_m,\ D>0\ ?\ \lfloor O.f\cdot D\rfloor : 256\,\mathrm{MB}\big)$\;
  }
  $(\mathit{modo},E,P)\gets(\mathsf{O1},0,1)$\;
  \lIf{$n\le\Bsz_0$}{\KwRet $(\mathit{modo},E,P,\beta_m)$}
  \eIf{$ns\le\beta_m$}{$(\mathit{modo},E)\gets(\mathsf{Completo},\,n)$\;}{
    $b\gets\lfloor\beta_m/s\rfloor$;\quad $E'\gets$ mayor potencia de 2 que no supera $b$\;
    \lIf{$E'>\Bsz_0$ \KwAnd $b\ge 2\Bsz_0$}{$(\mathit{modo},E)\gets(\mathsf{Bloques},\,E')$}
  }
  \If{$\mathit{modo}=\mathsf{Completo}$ \KwAnd $O.\mathit{hilos}\ne 1$}{
    $h\gets O.\mathit{hilos}$, o los hilos del hardware si es $0$\;
    $P\gets$ mayor potencia de 2 con $P\le h$ y $P\le\max(1,\lfloor n/2^{17}\rfloor)$\;
  }
  \KwRet $(\mathit{modo},E,P,\beta_m)$\;
}
\BlankLine
\Fn{\Rapido{$A, n, O$}}{
  $\pi\gets$ \PlanF{$n,\ \mathit{sizeof}(T),\ O$}\;
  \lIf{$n<2$}{\KwRet}
  semilla de \Aleatorio${}\gets\sigma\oplus n$, salvo que $O$ pida pivotes aleatorios\;
  \lIf{\Ordenado{$0,n$}}{\KwRet}
  \lIf{$T$ entero, ${\lt}={<}$, $O.\mathit{atajos}$ \KwAnd \Conteo{$A, n, \pi.\beta_m$}}{\KwRet}
  \While(\tcp*[f]{caída escalonada}){$\pi.\mathit{modo}\ne\mathsf{O1}$ \KwAnd \KwNot \Reservar{$\mathit{aux},\ \pi.E$}}{
    $e\gets(\pi.\mathit{modo}=\mathsf{Completo}\ ?\ \lfloor n/2\rfloor : \pi.E/2)$;\quad $E'\gets$ mayor potencia de 2 que no supera $e$;\quad $\pi.P\gets 1$\;
    \leIf{$E'>\Bsz_0$}{$(\pi.\mathit{modo},\pi.E)\gets(\mathsf{Bloques},E')$}{$(\pi.\mathit{modo},\pi.E)\gets(\mathsf{O1},0)$}
  }
  \lIf{$\pi.\mathit{modo}=\mathsf{Completo}$ \KwAnd $\pi.P>1$}{\Paralelo{$A, n, \mathit{aux}, \pi.P$}; \KwRet}
  \lIf{$\pi.\mathit{modo}\ne\mathsf{O1}$}{$\buf\gets\mathit{aux}$;\ $\Bsz\gets\pi.E$;\ $\mathit{galope}\gets\mathbf{true}$}
  \SotSort{$A, n$}\tcp*{el Algoritmo~\ref{alg:main}, con ese $\buf$ y ese $\Bsz$}
}
\end{algorithm}

\noindent $\sigma$ es una constante fija. \Reservar pide memoria sin lanzar
excepciones: devuelve falso si no la obtiene. Los buffers de 4\,MB o más se piden
alineados a 2\,MB y con \texttt{MADV\_HUGEPAGE} (páginas grandes, si el núcleo las da).

\begin{algorithm}[H]
\caption{Mezcla con buffer y galope (v8)}\label{alg:galope}
\Fn(\tcp*[h]{$m-l\le\Bsz$;\ $M_G=7$}){\MIzq{$l, m, r$}}{
  $a\gets m-l$;\quad $\buf[0..a)\gets A[l..m)$;\quad $p\gets 0$;\ $q\gets m$;\ $o\gets l$\;
  \lIf{\KwNot $\mathit{galope}$}{lazo clásico: un paso de cabeza por comparación; \KwRet}
  \While{$p<a$ \KwAnd $q<r$}{
    $c_x\gets 0$;\quad $c_y\gets 0$\tcp*{victorias seguidas de cada lado}
    \While(\tcp*[f]{uno a uno}){$p<a$ \KwAnd $q<r$ \KwAnd $\max(c_x,c_y)<M_G$}{
      \leIf{$A[q]\lt\buf[p]$}{$A[o{+}{+}]\gets A[q{+}{+}]$; $c_y\gets c_y{+}1$; $c_x\gets 0$}{$A[o{+}{+}]\gets\buf[p{+}{+}]$; $c_x\gets c_x{+}1$; $c_y\gets 0$}
    }
    \While(\tcp*[f]{galope}){$p<a$ \KwAnd $q<r$}{
      $k\gets$ \Galope{$\buf[p..a),\ x\mapsto A[q]\nlt x$}\tcp*{lo que no supera a la cabeza de $Y$}
      $A[o..o{+}k)\gets\buf[p..p{+}k)$;\quad $o\gets o{+}k$;\quad $p\gets p{+}k$\;
      \lIf{$p=a$}{\KwBreak}
      $k'\gets$ \Galope{$A[q..r),\ y\mapsto y\lt\buf[p]$}\tcp*{lo estrictamente menor que la cabeza de $X$}
      $A[o..o{+}k')\gets A[q..q{+}k')$;\quad $o\gets o{+}k'$;\quad $q\gets q{+}k'$\tcp*{$o<q$: hacia adelante}
      \lIf{$k<M_G$ \KwAnd $k'<M_G$}{\KwBreak}
    }
  }
  $A[o..o{+}a{-}p)\gets\buf[p..a)$\;
}
\BlankLine
\Fn(\tcp*[h]{$r-m\le\Bsz$}){\MDer{$l, m, r$}}{
  simétrico: la derecha va al buffer, se mezcla desde el final y \Galope busca desde el final\;
}
\BlankLine
\Fn(\tcp*[h]{$Q$ monótono: verdadero y después falso}){\Galope{$S[0..m),\ Q$}}{
  \lIf{$m=0$ \KwOr \KwNot $Q(S[0])$}{\KwRet $0$}
  $i\gets 0$;\quad $\delta\gets 1$\;
  \While(\tcp*[f]{prueba $S[1], S[3], S[7], S[15],\dots$}){$i+\delta<m$ \KwAnd $Q(S[i+\delta])$}{$i\gets i+\delta$;\quad $\delta\gets 2\delta$\;}
  $h\gets\min(i+\delta,\ m)$\;
  \KwRet el menor $j\in(i,h]$ con $j=h$ o $\neg Q(S[j])$, por búsqueda binaria\;
}
\end{algorithm}

\begin{algorithm}[H]
\caption{Conteo de enteros y modo paralelo con \emph{merge path} (v8)}\label{alg:paralelo}
\Fn(\tcp*[h]{$T$ entero y ${\lt}={<}$}){\Conteo{$A, n, \beta_m$}}{
  \lIf{$n<4096$}{\KwRet $\mathbf{false}$}
  \lIf{el rango de $A[0], A[s], A[2s],\dots$ con $s=\lfloor n/1024\rfloor$ es $\ge 2^{16}$}{\KwRet $\mathbf{false}$}
  $\mu\gets\min A$;\quad $R\gets\max A-\mu$\;
  \lIf{$R\ge 2^{16}$ \KwOr $8(R+1)>\beta_m$}{\KwRet $\mathbf{false}$}
  $c[0..R]\gets 0$;\quad \lFor{$i\gets 0$ \KwTo $n-1$}{$c[A[i]-\mu]\gets c[A[i]-\mu]+1$}
  $o\gets 0$;\quad \lFor{$v\gets 0$ \KwTo $R$}{$A[o..o{+}c[v])\gets\mu+v$;\ $o\gets o+c[v]$}
  \KwRet $\mathbf{true}$\;
}
\BlankLine
\Fn(\tcp*[h]{cuántos de $X$ hay entre los primeros $k$ de $\Mrg(X,Y)$}){\CoRango{$k,\ X[0..a),\ Y[0..b)$}}{
  $lo\gets\max(0,k-b)$;\quad $hi\gets\min(k,a)$\;
  \While{$lo<hi$}{
    $i\gets\lfloor(lo+hi)/2\rfloor$;\quad $j\gets k-i$\;
    \leIf(\tcp*[f]{$X[i]$ va antes que $Y[j-1]$}){$Y[j-1]\nlt X[i]$}{$lo\gets i+1$}{$hi\gets i$}
  }
  \KwRet $lo$\;
}
\BlankLine
\Fn(\tcp*[h]{$P$ potencia de 2; $\mathit{aux}$ de $n$ elementos}){\Paralelo{$A, n, \mathit{aux}, P$}}{
  $d_i\gets\lfloor ni/P\rfloor$ para $i=0,\dots,P$\;
  \textbf{en paralelo} para $i<P$: \SotSort sobre $A[d_i..d_{i+1})$ con $\buf=\mathit{aux}[d_i..d_{i+1})$ y galope\;
  $S\gets A$;\quad $D\gets\mathit{aux}$;\quad $w\gets 1$\;
  \For{$\rho\gets P, P/2,\dots,2$}{
    $g\gets 2P/\rho$\tcp*{tareas por cada par de runs}
    \ForEach(\tcp*[f]{par $u=\lfloor t/g\rfloor$, parte $v=t\bmod g$}){tarea $t<P$}{
      $lo\gets d_{2uw}$;\ $md\gets d_{(2u+1)w}$;\ $hi\gets d_{(2u+2)w}$;\ $N\gets hi-lo$\;
      $k_0\gets\lfloor Nv/g\rfloor$;\ $k_1\gets\lfloor N(v{+}1)/g\rfloor$;\ $i_\ell\gets$ \CoRango{$k_\ell,\ S[lo..md),\ S[md..hi)$} para $\ell=0,1$\;
    }
    \textbf{en paralelo} para cada tarea: mezclar $S[lo{+}i_0..lo{+}i_1)$ con $S[md{+}k_0{-}i_0..md{+}k_1{-}i_1)$ en $D[lo{+}k_0..lo{+}k_1)$\;
    intercambiar $S$ y $D$;\quad $w\gets 2w$\;
  }
  \lIf{$S\ne A$}{\textbf{en paralelo}, copiar $S$ a $A$ por trozos}
}
\end{algorithm}

\noindent Los cortes de cada ronda se calculan todos antes de mezclar: con tipos que
no se copian byte a byte, mover vacía el origen, y otro hilo podría estar buscando su
corte sobre esos mismos elementos (ThreadSanitizer encontró esa carrera en una
versión anterior del código).

\begin{algorithm}[H]
\caption{Orden estable por clave entera (v8)}\label{alg:clave}
\KwIn{$A[0..n)$ y una clave entera $\kappa$; el orden es $x\lt_\kappa y\iff\kappa(x)<\kappa(y)$}
\Fn{\PorClave{$A, n, \kappa, O$}}{
  $\pi\gets$ \PlanF{$n,\ \mathit{sizeof}(T),\ O$}\;
  \If{$T$ no se copia byte a byte \KwOr $n<4096$ \KwOr $\pi.\mathit{modo}\ne\mathsf{Completo}$}{\Rapido{$A, n, O$} con $\lt_\kappa$;\quad \KwRet\;}
  \lIf{$\kappa(A[0]),\dots,\kappa(A[n-1])$ es no decreciente}{\KwRet}
  \lIf{\PocosGiros{$n$}}{\Rapido{$A, n, O$} con $\lt_\kappa$; \KwRet}
  $\mu\gets\min_i\kappa(A[i])$;\quad $R\gets\max_i\kappa(A[i])-\mu$;\quad $r\gets\mathrm{bits}(R)$;\quad $\iota\gets\mathrm{bits}(n-1)$\;
  \If(\tcp*[f]{indirecto empaquetado}){$\mathit{sizeof}(T)\ge 32$ \KwAnd $r+\iota\le 64$ \KwAnd $n(\mathit{sizeof}(T)+16)\le\pi.\beta_m$}{
    $K[i]\gets(\kappa(A[i])-\mu)\cdot 2^{\iota}+i$ para $i<n$\tcp*{enteros de 64 bits distintos}
    \Rapido{$K, n$} con $<$, sin sondeo y con tope de $8n$ bytes\;
    $\mathit{aux}[j]\gets A[K[j]\bmod 2^{\iota}]$ para $j<n$;\quad $A\gets\mathit{aux}$;\quad \KwRet\;
  }
  \If(\tcp*[f]{radix LSD sobre los registros}){$R<2^{24}$}{
    $q\gets 1$, $2$ o $3$ según $R<2^8$, $2^{16}$ o $2^{24}$;\quad $S\gets A$;\quad $D\gets\mathit{aux}$\;
    \For{$\delta\gets 0$ \KwTo $q-1$}{distribuir $S$ en $D$, en orden, por el byte $\delta$ de $\kappa(x)-\mu$;\quad intercambiar $S$ y $D$\;}
    \lIf{$S\ne A$}{$A\gets S$}
    \KwRet\;
  }
  \Rapido{$A, n, O$} con $\lt_\kappa$\;
}
\end{algorithm}

\noindent $\mathrm{bits}(x)$ es la cantidad de bits de $x$ ($\mathrm{bits}(0)=0$). Si una
reserva falla, se pasa a la rama siguiente.


\subsection{Mezcla con galope}

\begin{lema}[Galope]\label{lem:galope}
Sea $Q$ monótono sobre $S[0..m)$: existe $k^*$ con $Q(S[i])\iff i<k^*$. Entonces
\Galope{$S,Q$} devuelve $k^*$ y evalúa $Q$ a lo sumo $2\lfloor\log_2 k^*\rfloor+2$ veces
si $k^*\ge1$, y una vez si $k^*=0$ (ninguna si $m=0$). En los dos casos son a lo sumo
$2\lfloor\log_2(k^*{+}1)\rfloor+2$.
\end{lema}
\begin{proof}
Si $m=0$ o $\neg Q(S[0])$, $k^*=0$. Si no, el lazo exponencial mantiene $Q(S[i])$, es
decir, $i<k^*$; la variante es $m-i$, porque $i$ crece en cada vuelta. Al salir, o bien
$i+\delta\ge m$ y entonces $k^*\le m=h$, o bien $\neg Q(S[i+\delta])$ y entonces
$k^*\le i+\delta=h$. En los dos casos $k^*\in(i,h]$. Para $j\in(i,k^*)$ vale $Q(S[j])$, y
en $j=k^*$ vale $j=h$ o $\neg Q(S[j])$: el menor $j$ de la línea final es $k^*$, y la
búsqueda binaria lo encuentra porque la condición es monótona en $j$.

\emph{Costo.} Las posiciones probadas por el lazo son $2^{t+1}-1$ para $t=0,1,\dots$;
las que dan verdadero son las menores que $k^*$, o sea las $t$ con
$2^{t+1}\le k^*$: a lo sumo $\lfloor\log_2k^*\rfloor$. Hay a lo sumo una prueba que da
falso. El intervalo que queda para la búsqueda binaria tiene a lo sumo
$\delta-1=2^{\lfloor\log_2 k^*\rfloor}-1$ posiciones, que se resuelven con a lo sumo
$\lfloor\log_2k^*\rfloor$ evaluaciones. Con la evaluación de $S[0]$, el total es a lo
sumo $2\lfloor\log_2k^*\rfloor+2$. La búsqueda desde el final de \MDer es simétrica,
con posiciones $1, 2, 4, 8,\dots$ contadas desde el final, y tiene la misma cota.
\end{proof}

\begin{lema}[Monotonía]\label{lem:galmono}
Si $X$ está ordenada, $x\mapsto y\nlt x$ es monótono sobre $X$ para todo $y$. Si $Y$
está ordenada, $y\mapsto y\lt x$ es monótono sobre $Y$ para todo $x$.
\end{lema}
\begin{proof}
Sean $j<i$. Como $X$ está ordenada, $X[j]\les X[i]$. Si $y\nlt X[i]$, es decir,
$X[i]\les y$, la transitividad de $\les$ (Lema~\ref{lem:swo}(d)) da $X[j]\les y$. En
el segundo caso, si $Y[i]\lt x$, el Lema~\ref{lem:swo}(c) con $Y[j]\les Y[i]$ da
$Y[j]\lt x$.
\end{proof}

\begin{lema}[MergeIzq con galope]\label{lem:mizqg}
Con $\mathit{galope}$ verdadero, \MIzq cumple la misma tripla del
Lema~\ref{lem:mizq}, con a lo sumo $6(|X|+|Y|)$ comparaciones. Lo mismo vale para
\MDer.
\end{lema}
\begin{proof}
Se usa el invariante del Lema~\ref{lem:mizq}: $o=l+p+(q-m)$,
$A[l..o)=\Mrg(X[0..p),Y[0..q{-}m))$, $\buf[p..a)=X[p..)$ y $A[q..r)=Y[q{-}m..)$. Como
$o-q=p-a$, vale $o\le q$, con desigualdad estricta mientras $p<a$.

Cada paso del lazo uno a uno es un paso de cabeza (Lema~\ref{lem:merge}(a)). En el
galope sobre $X$, la cabeza de $Y$ es fija, $y=A[q]$, y por los
Lemas~\ref{lem:galmono} y~\ref{lem:galope} los $k$ elementos copiados son exactamente
los primeros de $X[p..)$ que cumplen $y\nlt x$. Para cada uno, el paso de cabeza elige
$X$, porque no vale $Y[\cdot]\lt X[\cdot]$; así, la copia equivale a $k$ pasos de
cabeza. Si después $p<a$, el elemento $X[p]$ no cumple el predicado: $A[q]\lt X[p]$. En
el galope sobre $Y$ la cabeza de $X$ es fija, $x=\buf[p]$, y los $k'$ elementos
copiados son los primeros de $Y[q{-}m..)$ estrictamente menores que $x$: son $k'$ pasos
de cabeza que eligen $Y$. Como $A[q]\lt\buf[p]$, vale $k'\ge 1$. La copia de $Y$ va
hacia adelante con $o<q$, así que no pisa nada sin leer.

\emph{Terminación.} Cada vuelta del lazo externo hace al menos un paso uno a uno, y
cada vuelta del galope o deja $p=a$ o copia $k'\ge1$ elementos: la variante
$(a-p)+(r-q)$ decrece. Al salir, el final es el del Lema~\ref{lem:mizq}.

\emph{Comparaciones.} En el lazo uno a uno, una por elemento escrito. Una vuelta del
galope que hace las dos búsquedas cuesta, por el Lema~\ref{lem:galope}, a lo sumo
$2\lfloor\log_2(k{+}1)\rfloor+2+2\lfloor\log_2(k'{+}1)\rfloor+2\le 2k+2k'+4\le 6(k+k')$,
porque $k'\ge1$. Una vuelta que termina con $p=a$ copió $k\ge1$ elementos y cuesta a lo
sumo $2k+2\le 4k$. Cada elemento se escribe una sola vez, así que el total es a lo sumo
$6(|X|+|Y|)$. \MDer es simétrico, con los pasos de cola del Lema~\ref{lem:merge}(b).
\end{proof}

\begin{observacion}
Con galope, en el Lema~\ref{lem:sym}(d) las mezclas de las hojas comparan a lo sumo 6
veces por elemento en lugar de una: la cota pasa a
$6s+\bigl(1+\frac{2s}{\Bsz}\bigr)(5\ell+2)$. Las cotas asintóticas del
Lema~\ref{lem:respaldo}, del Lema~\ref{lem:natcost} y del Teorema~\ref{thm:cost} no
cambian; solo sus constantes. El galope está apagado en el modo $O(1)$ para tipos
aritméticos, que es como se midió y demostró v7.
\end{observacion}

\subsection{Conteo de enteros}

\begin{lema}[Conteo]\label{lem:conteo}
Bajo el Supuesto~\ref{sup:enteros},
\[
\hoare{A=A_0}{c\gets\Conteo(A,n,\beta_m)}{(\neg c\wedge A=A_0)\ \vee\ (c\wedge A=\Omega)},
\]
con $O(n+R)$ operaciones y $8(R+1)\le\beta_m$ bytes de contadores, que se liberan antes
de volver.
\end{lema}
\begin{proof}
Hasta la línea 5, \Conteo solo lee. La muestra es un subconjunto de $A$, así que su
rango no supera a $R$: descartar por la muestra es descartar un caso con $R\ge2^{16}$.
Si devuelve falso, $A=A_0$.

Las diferencias $A[i]-\mu$ se calculan en aritmética sin signo de $w$ bits, módulo
$2^w$. Como $0\le A[i]-\mu\le R\le 2^w-1$, el resultado es la diferencia verdadera y el
índice está en $[0,R]$.

Invariante del primer lazo: tras $i$ vueltas, $c[v]=\#\{j<i:\ A_0[j]=\mu+v\}$. Invariante
del segundo: tras procesar los valores $0,\dots,v-1$, vale $o=\sum_{u<v}c[u]$ y $A[0..o)$
es la secuencia no decreciente que repite cada $\mu+u$, para $u<v$, exactamente $c[u]$
veces. Al final $o=n$ y $A$ es la única secuencia no decreciente con el multiconjunto de
valores de $A_0$. $\Omega$ está ordenada, que para enteros con $<$ significa no
decreciente, y tiene ese mismo multiconjunto: las dos secuencias de valores coinciden.
Por el Supuesto~\ref{sup:enteros}, $A=\Omega$.
\end{proof}

\subsection{Modo paralelo}

\begin{lema}[CoRango]\label{lem:corango}
Sean $X[0..a)$ e $Y[0..b)$ ordenadas, $Z=\Mrg(X,Y)$, $0\le k\le a+b$, y sea $c(k)$ la
cantidad de elementos de $X$ en $Z[0..k)$. Entonces \CoRango{$k,X,Y$} devuelve $c(k)$
con a lo sumo $\lceil\log_2(\min(k,a)-\max(0,k-b)+1)\rceil$ comparaciones, $c$ es no
decreciente, y para $k_0\le k_1$:
\[
Z[k_0..k_1)=\Mrg\big(X[c(k_0)..c(k_1)),\ Y[k_0{-}c(k_0)..k_1{-}c(k_1))\big).
\]
\end{lema}
\begin{proof}
Como $\Mrg$ preserva el orden interno de $X$ y de $Y$, los elementos de $X$ en
$Z[0..k)$ son $X[0..c)$ y los de $Y$ son $Y[0..k{-}c)$, con $c=c(k)$. Por eso
$\max(0,k-b)\le c\le\min(k,a)$.

\emph{Afirmación.} Para $i\in[\max(0,k-b),\min(k,a))$ y $j=k-i$ (vale $1\le j\le b$):
$i<c\iff Y[j-1]\nlt X[i]$.
($\Rightarrow$) Si $i<c$, $X[i]$ está en $Z[0..k)$. Como $k-c<k-i=j$, $Y[j-1]$ no está
en $Z[0..k)$, así que $X[i]$ precede a $Y[j-1]$ en $Z$, que está ordenada:
$Y[j-1]\nlt X[i]$.
($\Leftarrow$) Si $i\ge c$, $X[i]$ no está en $Z[0..k)$ y, como $k-c\ge j$, $Y[j-1]$
sí. Entonces $Y[j-1]$ precede a $X[i]$ y, por estar ordenada, $Y[j-1]\les X[i]$. Si
fueran equivalentes, la Definición~\ref{def:merge} pondría a $X[i]$ primero; luego
$Y[j-1]\lt X[i]$.

La búsqueda binaria mantiene $lo\le c\le hi$: si la condición vale, $i<c$ y
$lo\gets i+1\le c$; si no, $i\ge c$ y $hi\gets i$. Los índices son válidos porque
$lo\le i<hi$ da $0\le i<a$ y $1\le j\le b$. La variante es $hi-lo$, que al menos se
reduce a la mitad, lo que da la cota de comparaciones. Al terminar, $lo=hi=c$.

$c$ es no decreciente porque $Z[0..k)$ crece con $k$. La igualdad final es el
Lema~\ref{lem:merge}(c): $Z[k_0..k_1)$ es un segmento de $Z$ y contiene exactamente esos
elementos de $X$ y de $Y$.
\end{proof}

\begin{lema}[Paralelo]\label{lem:paralelo}
Bajo los Supuestos~\ref{sup:hilos} y~\ref{sup:reserva}, con $P=2^p\le n$ y
$\mathit{aux}$ de $n$ celdas:
\[
\hoare{A=A_0}{\Paralelo(A,n,\mathit{aux},P)}{A=\Omega}.
\]
Además, ninguna celda es escrita por dos hilos en la misma fase, y ninguna celda que un
hilo lee en una fase es escrita por otro en esa fase.
\end{lema}
\begin{proof}
\emph{Fase de trozos.} El hilo $i$ solo toca $A[d_i..d_{i+1})$ y
$\mathit{aux}[d_i..d_{i+1})$, que son disjuntos entre hilos. Las llamadas a $\lt$ son
seguras por el Supuesto~\ref{sup:hilos}. Por el Teorema~\ref{thm:main}, con
$\Bsz=d_{i+1}-d_i$ (Observación~\ref{obs:bsz}) o con el buffer propio si el trozo es más
chico, cada trozo queda ordenado y la transformación es OE dentro del trozo. Por el
Lema~\ref{lem:oe}(c), es OE sobre todo el arreglo, que queda coherente.

\emph{Invariante de rondas.} Al empezar la ronda de ancho $w$, con $\rho=P/w$ runs:
$S$ contiene cada elemento de $A_0$ una vez, es coherente, y cada run
$S[d_{jw}..d_{(j+1)w})$ está ordenado. Sea un par $X=S[lo..md)$, $Y=S[md..hi)$ y
$Z=\Mrg(X,Y)$. Sus $g$ tareas usan los cortes $\lfloor Nv/g\rfloor$, $v=0,\dots,g$, que
parten $[0,N)$ en tramos consecutivos. Por el Lema~\ref{lem:corango}, la tarea $v$
escribe $D[lo{+}k_0..lo{+}k_1)=Z[k_0..k_1)$, y entre todas escriben
$D[lo..hi)=Z$. $\Mrg$ es OE (Lema~\ref{lem:merge}(d)) y los pares son disjuntos, así que
$D$ es coherente y sus runs de ancho $2w$ están ordenados. Al intercambiar, el invariante
vale para $2w$.

\emph{Carreras.} En el cálculo de cortes solo se lee $S$. En la mezcla, cada tarea
escribe un tramo de $D$ disjunto de los demás y solo lee $S$, que nadie escribe en esa
fase: $A$ y $\mathit{aux}$ son disjuntos (Supuesto~\ref{sup:reserva}). Las barreras
separan las fases (Supuesto~\ref{sup:hilos}). La copia final escribe tramos disjuntos de
$A$ y lee solo $S$. Que los cortes se calculen todos antes de mezclar es necesario:
para tipos no triviales, mover deja el origen en un estado no especificado, y una
búsqueda de corte concurrente leería esos valores.

\emph{Final.} Tras $p$ rondas queda un solo run: $S[0..n)$ está ordenado y es coherente,
así que $S=\Omega$ por el Lema~\ref{lem:coh}. Si $S\ne A$, la copia lo lleva a $A$
preservando identidades (Supuesto~\ref{sup:ident}). Termina porque hay $p$ rondas y cada
lazo es finito.
\end{proof}

\subsection{Corrección, memoria y cotas del modo rápido}

\begin{teorema}[Corrección de \Rapido]\label{thm:rapido}
Bajo los Supuestos~\ref{sup:orden}--\ref{sup:n} y~\ref{sup:reserva}, más el
Supuesto~\ref{sup:hilos} si $P>1$ y el Supuesto~\ref{sup:enteros} si se usa \Conteo:
\[
\hoare{A=A_0,\ |A_0|=n}{\Rapido(A,n,O)}{A=\Omega}.
\]
\end{teorema}
\begin{proof}
\PlanF solo calcula números y no toca $A$. Si $n<2$ no hay nada que hacer. Si \Ordenado
devuelve verdadero, $A$ queda ordenado y OE (Lema~\ref{lem:yaord}) y, por el
Lema~\ref{lem:coh}, $A=\Omega$. Si \Conteo devuelve verdadero, $A=\Omega$ por el
Lema~\ref{lem:conteo}; si devuelve falso, $A=A_0$. El lazo de reservas termina porque
el modo baja estrictamente ($\mathsf{Completo}\to\mathsf{Bloques}$ con $E$ menor o
$\mathsf{O1}$, y en $\mathsf{Bloques}$ $E$ se reduce a la mitad) y en $\mathsf{O1}$ se
detiene; no modifica $A$. Con $P>1$, $P\le n/2^{17}$ por \PlanF y el
Lema~\ref{lem:paralelo} da $A=\Omega$. En otro caso se ejecuta \SotSort con $\Bsz=E$ o
$\Bsz=\Bsz_0$, y con galope si hay buffer: el Teorema~\ref{thm:main} vale por la
Observación~\ref{obs:bsz} y el Lema~\ref{lem:mizqg}. \SotSort vuelve a llamar a
\Ordenado, que ya devolvió falso sobre el mismo $A$ y vuelve a hacerlo sin modificarlo
(Lema~\ref{lem:yaord}).
\end{proof}

\begin{teorema}[Memoria de \Rapido]\label{thm:rapmem}
Además del arreglo, \Rapido usa:
\begin{enumerate}[label=(\alph*),nosep]
\item a lo sumo $\beta_m$ bytes reservados en todo momento: el buffer
  $\mathit{aux}$, de $E\cdot\mathit{sizeof}(T)\le\beta_m$ bytes, o los contadores de
  \Conteo, de $8(R+1)\le\beta_m$ bytes, nunca los dos a la vez;
\item la memoria fija del Teorema~\ref{thm:mem} (buffer propio y pilas) por cada hilo
  que ordena un trozo;
\item $O(P)$ palabras para los bordes, los cortes y los hilos, más la pila que el
  sistema le asigna a cada uno de los $P-1$ hilos adicionales.
\end{enumerate}
\end{teorema}
\begin{proof}
(a) En modo completo, $ns\le\beta_m$. En modo bloques, $E\le b=\lfloor\beta_m/s\rfloor$,
así que $Es\le\beta_m$. En la caída escalonada, desde completo
$E'\le\lfloor n/2\rfloor$ y desde bloques $E'\le E/2$: el tamaño solo baja. \Conteo
verifica $8(R+1)\le\beta_m$ y libera los contadores antes de volver, y la reserva del
buffer ocurre después. (b) y (c) son inspección del Algoritmo~\ref{alg:paralelo}.
\end{proof}

\noindent El presupuesto es $\beta_m=\min(O.\mathit{max},\ O.f\cdot D)$, con $D$ la
memoria disponible \emph{en el momento de la consulta}. El teorema garantiza que no se
pide más que eso; no garantiza que otros procesos no consuman la memoria restante
mientras se ordena.

\begin{teorema}[Cotas con buffer completo]\label{thm:rapcost}
Si \SotSort se ejecuta con $\Bsz\ge n$, hace $O(n\log n)$ comparaciones y $O(n\log n)$
movimientos en el peor caso, con galope o sin él. \Paralelo con $\mathit{aux}$ de $n$
celdas hace en total $O(n\log n)$ comparaciones y movimientos.
\end{teorema}
\begin{proof}
\emph{Camino \Quick.} Por el Lema~\ref{lem:prof}(b), los intervalos procesados que
contienen a un elemento forman una cadena de longitud $O(\log n)$. En un intervalo de
tamaño $\lambda\le n\le\Bsz$, cada \Part tiene $k=\lceil\lambda/\Bsz\rceil=1$ bloque: por
el Lema~\ref{lem:partest}, $\lambda$ comparaciones, a lo sumo $3\lambda$ escrituras y
ninguna rotación. Cada iteración hace a lo sumo dos particiones y elige un pivote con
$O(\lambda^{1/3}\log\lambda)$ operaciones. Los intervalos de una misma profundidad son
disjuntos, así que la suma es $O(n\log n)$. Las hojas son disjuntas y cuestan
$O(\leaf\log\leaf)$ cada una (Lema~\ref{lem:hoja}).

\emph{Respaldo y SymMerge.} En todo SymMerge de tamaño $s\le n\le\Bsz$ vale
$\min(n_l,n_r)\le\Bsz$, así que la raíz es una hoja y no hay división: a lo sumo $3s$
escrituras (Lema~\ref{lem:sym}(c)) y $6s+2\lceil\log_2(s{+}1)\rceil+1$ comparaciones
(Lema~\ref{lem:mizqg} y las búsquedas del recorte). \Respaldo sobre un intervalo de
tamaño $\lambda$ hace $O(\lambda)$ por nivel y $O(\log\lambda)$ niveles, más la inserción
en bloques de 24; los intervalos de respaldo son disjuntos.

\emph{Camino natural.} Por el Lema~\ref{lem:natcost}(b), cada elemento participa en
$O(\log n)$ fusiones, cada una lineal por lo anterior; la detección y la inserción de
runs cortos son $O(n)$.

\emph{Paralelo.} Los trozos cuestan $O(\frac{n}{P}\log\frac{n}{P})$ cada uno por lo
anterior. Cada una de las $\log_2P$ rondas hace $n$ movimientos, a lo sumo $n$
comparaciones de mezcla y $2P$ llamadas a \CoRango de $O(\log n)$ comparaciones. Como
$P\le n$, el total es $O(n\log n)$.
\end{proof}

\noindent En modo bloques valen las cotas del Teorema~\ref{thm:cost} con $\Bsz=E$:
$O(n\log n\log(n/E))$ movimientos y $O(n\log n+\frac{n}{E}\log^2n)$ comparaciones.

\begin{teorema}[Determinismo]\label{thm:det}
\begin{enumerate}[label=(\alph*),nosep]
\item El resultado es $\Omega$ cualquiera sea el plan, la semilla, la cantidad de hilos
  y el resultado de las reservas.
\item Si $O.\mathit{sondear}$ es falso, el plan depende solo de $n$,
  $\mathit{sizeof}(T)$, $\Bsz_0$ y $O$, y además de la cantidad de hilos del hardware si
  $O.\mathit{hilos}=0$.
\item Fijado el plan, si ninguna reserva falla y los pivotes no son aleatorios, la
  secuencia de operaciones de cada hilo es función de la entrada. Solo varía el
  entrelazado en el tiempo entre hilos, que no cambia ningún valor leído.
\end{enumerate}
\end{teorema}
\begin{proof}
(a) Es el Teorema~\ref{thm:rapido}. (b) \PlanF solo llama a \Disp si
$O.\mathit{sondear}$ es verdadero, y solo consulta el hardware si $O.\mathit{hilos}=0$.
(c) Cada decisión depende de $A$, del plan y del generador, cuya semilla es
$\sigma\oplus n$ (y una semilla fija por trozo). En \Paralelo, por el
Lema~\ref{lem:paralelo}, cada tarea lee solo valores que no cambian durante su fase.
\end{proof}

\subsection{Orden por clave}

Sea $\kappa$ una clave entera y $x\lt_\kappa y\iff\kappa(x)<\kappa(y)$, que es un orden
débil estricto porque $<$ lo es sobre los enteros. $\Omega$ se toma respecto de
$\lt_\kappa$, con etiquetas iguales a los índices.

\begin{lema}[Indirecto empaquetado]\label{lem:indirecto}
Sean $d_i=\kappa(A_0[i])-\mu\in[0,R]$, $r=\mathrm{bits}(R)$, $\iota=\mathrm{bits}(n-1)$
con $r+\iota\le 64$, y $K[i]=d_i\cdot2^\iota+i$. Entonces los $K[i]$ son distintos,
menores que $2^{64}$, y $K[i]<K[j]\iff A_0[i]\lts A_0[j]$. En consecuencia, si $K$ queda
ordenado de forma creciente, $A[j]\gets A_0[K[j]\bmod 2^\iota]$ deja $A=\Omega$.
\end{lema}
\begin{proof}
$i<2^\iota$ y $d_i<2^r$, así que $K[i]<2^{r+\iota}\le2^{64}$, y $K[i]\bmod2^\iota=i$,
$\lfloor K[i]/2^\iota\rfloor=d_i$: la codificación es inyectiva. Como $i,j<2^\iota$,
comparar $K$ es comparar lexicográficamente $(d,\text{índice})$:
$K[i]<K[j]\iff d_i<d_j\vee(d_i=d_j\wedge i<j)$. Además $d_i<d_j\iff A_0[i]\lt_\kappa
A_0[j]$ y $d_i=d_j\iff A_0[i]\eqv A_0[j]$, que con etiquetas iguales a los índices es la
Definición~\ref{def:omega} de $\lts$. Las restas se hacen sin signo, módulo $2^w$, y son
exactas porque $0\le\kappa-\mu\le R<2^w$. $K$ tiene valores distintos, así que su orden
creciente es único; el recorrido final pone los elementos en orden creciente según
$\lts$, que es $\Omega$. $\mathit{aux}$ es disjunto de $A$ y la copia preserva
identidades.
\end{proof}

\begin{lema}[Radix por clave]\label{lem:radix}
Sea $q$ tal que $R<2^{8q}$ y $\kappa_\delta(x)=(\kappa(x)-\mu)\bmod 2^{8(\delta+1)}$.
Tras la pasada $\delta$ (contando desde 0), $S$ contiene cada elemento una vez, en orden
creciente según el par $(\kappa_\delta(x),\etq(x))$. Al final, $A=\Omega$.
\end{lema}
\begin{proof}
Antes de la primera pasada, $S=A_0$ está en orden creciente de etiqueta. Una pasada
reparte por el byte $\delta$ con sumas de prefijos y un recorrido en orden de $S$: los
elementos quedan en orden creciente de byte y, a igual byte, en el orden que tenían en
$S$. Como $\kappa_\delta=\text{byte}_\delta\cdot2^{8\delta}+\kappa_{\delta-1}$, el
nuevo orden es creciente según $(\text{byte}_\delta,\kappa_{\delta-1},\etq)$, es decir,
según $(\kappa_\delta,\etq)$. Tras $q$ pasadas, $\kappa_{q-1}=\kappa-\mu$ porque
$R<2^{8q}$, y ese orden es $\lts$. Si el resultado quedó en $\mathit{aux}$, se copia.
\end{proof}

\begin{teorema}[\PorClave]\label{thm:clave}
Bajo los mismos supuestos que el Teorema~\ref{thm:rapido},
\[\hoare{A=A_0}{\PorClave(A,n,\kappa,O)}{A=\Omega},\]
y la memoria reservada en todo
momento es a lo sumo $\beta_m$ bytes más $O(1)$ (los contadores del radix).
\end{teorema}
\begin{proof}
Las ramas que delegan en \Rapido con $\lt_\kappa$ son correctas por el
Teorema~\ref{thm:rapido}. Si las claves ya son no decrecientes, $A_0$ está ordenada y es
coherente: $A_0=\Omega$. La rama indirecta es correcta por el Lema~\ref{lem:indirecto}
y el Teorema~\ref{thm:rapido} aplicado a $K$; la de radix, por el
Lema~\ref{lem:radix}. Memoria: la rama indirecta tiene a la vez $K$ ($8n$ bytes),
$\mathit{aux}$ ($n\,\mathit{sizeof}(T)$) y lo que use \Rapido sobre $K$, acotado por su
tope de $8n$ (Teorema~\ref{thm:rapmem}); en total $n(\mathit{sizeof}(T)+16)\le\beta_m$,
que la rama verifica. La de radix usa $\mathit{aux}$, que entra porque el modo es
completo, y $3\cdot256$ contadores.
\end{proof}
